\ficha{Triner (1996), Banking, economic growth and industrialization: Brazil, 1906-30}{triner1996}

\begin{identificacao}
TRINER, Gail D. Banking, economic growth and industrialization: Brazil, 1906-30. \textbf{Revista Brasileira de Economia}, Rio de Janeiro, v. 50, n. 1, p. 135-153, jan./mar. 1996.

Artigo de periódico, 19 páginas, lido integralmente, em inglês, com resumo em português. Versão lida é a publicada. Deriva de capítulo da tese de doutorado da autora, Columbia University, 1994, que é a fonte das séries de balanços.
\end{identificacao}

\bloco{Problema e tese}

Que relação houve entre o sistema bancário brasileiro e o crescimento econômico e a industrialização de 1906 a 1930. A resposta é que os bancos se comportaram como preveem as teorias formuladas para economias desenvolvidas em fases equivalentes, integrados à economia produtiva e mais sensíveis à produção industrial do que à agrícola. O achado que a autora destaca é institucional: a função de política pública do Banco do Brasil obscurece nos agregados o dinamismo dos bancos privados (p. 136).

\bloco{Argumento}

\begin{enumerate}[leftmargin=*,nosep]
\item A reforma de 1906 encerra a estagnação que se seguiu às quebras de 1900 e 1901, quando a economia ficou praticamente sem sistema bancário organizado (p. 137). Seu traço decisivo é a abertura do Banco do Brasil, que acumula desde o início as funções de banco central do governo e de maior banco comercial privado (p. 137). Entre 1907 e 1930 ele deteve de 10 a 37\% dos depósitos bancários totais (p. 137, nota 8).

\item O crédito era pouco elástico por construção. Os bancos operavam a curto prazo com colateral identificável, e os estatutos vedavam empréstimo longo não garantido por imóvel, inclusive o de formação de capital industrial (p. 137). A renovação no vencimento convertia curto em médio prazo, mas os bancos não expandiam agressivamente seus passivos para criar crédito, de modo que a queda da razão caixa-depósito não gerou política creditícia mais expansiva (p. 137). Daí a formulação mais forte do texto: o sistema executava parcela crescente das transações sem criar recursos financeiros (p. 138).

\item A bancarização avançou depressa mesmo assim. Os depósitos passaram de 25\% de M2 em 1906 para 62\% em 1930 (p. 139), e os depósitos reais per capita se multiplicaram por 5,4, contra 2,1 da renda per capita (p. 140). A consequência monetária é explícita: a capacidade das autoridades de controlar a oferta de moeda apenas pelo volume de papel-moeda diminuiu (p. 139).

\item A política monetária era o instrumento com que o governo arbitrava concorrências, entre serviço da dívida externa, comércio do café e produção interna, mais do que a política fiscal, e o governo financiou grandes expansões de moeda e de crédito pelo Banco do Brasil (p. 138).

\item Nas regressões de 1907 a 1930, preços e produto explicam 39\% da variação dos depósitos reais no sistema total e cerca de 70\% nos bancos privados domésticos e nos paulistas (p. 148). Os depósitos respondem positiva e significativamente ao produto industrial, e não respondem na direção prevista ao agrícola (p. 148). O Banco do Brasil é a exceção sistemática, com $R^2$ de 0,19 e coeficientes de sinal invertido (p. 146, p. 149).
\end{enumerate}

\chave{In many ways, the Banco do Brasil functioned as a central bank for the federal government. It was also the largest commercial bank for the private sector.}{p. 137}

\chave{So, although the banking system executed an increasing share of the transactions undertaken by the economy, it did not create financial resources.}{p. 138}

\bloco{Conceitos e definições operacionais}

Elasticidade do crédito não é definida em abstrato, e sim por dois indicadores de balanço, a razão caixa sobre depósitos e a expansão deliberada de passivos (p. 137). Autoridade monetária é definida por função, não por estatuto, e a autora sustenta que tratar o Banco do Brasil como banco comercial comum distorce a análise da oferta de moeda, crítica dirigida a Neuhaus (p. 150, nota 36). Oferta de moeda é M2, com depósitos à vista e a prazo combinados por consistência de fonte (p. 136, nota 6). Demanda por depósitos reais é função de renda e de preço, com a variação de preços como proxy do juro na ausência de série (p. 143). O texto não define padrão-ouro, conversibilidade, lastro nem paridade.

\bloco{Evidência e método}

Balanços dos maiores bancos estrangeiros e domésticos, publicados em relatórios anuais e na imprensa financeira, agregados como representação do sistema e correspondentes a 70 a 85\% dos saldos bancários (p. 136). Fonte principal, o Jornal do Commercio, com O Estado de S. Paulo, o Anuário Estatístico do Estado de São Paulo e fontes regionais de Minas Gerais e Rio Grande do Sul (p. 136, nota 5). Produto e preços vêm de IBGE (1990), Haddad (1974) e Catão. Modelo de duas equações simultâneas para oferta e demanda de depósitos reais, por mínimos quadrados de dois estágios, variáveis em taxa anual de variação, N igual a 24, sem defasagens (p. 143-144). Não há contrafactual, e os testes de Granger-Sims sobre antecedência são indeterminados (p. 144, nota 27).

\bloco{Agentes e vertentes}

Schumpeter, Gerschenkron e Cameron sustentam que o sistema bancário influencia crescimento e industrialização; Robinson trata a expansão bancária como resíduo (p. 135). Na historiografia brasileira, Neuhaus e Peláez são apontados como os proponentes mais específicos da hipótese de economia subbancarizada, ao lado de Peláez e Suzigan e de Villela e Suzigan; Fritsch é citado para a tese de que o serviço da dívida externa comandava a política monetária, e Topik para o intervencionismo estatal no setor financeiro (p. 136, nota 4). Cano é contestado quanto à autonomia entre regiões (p. 149). Nenhum agente do debate público de 1906 a 1914 é nomeado.

\bloco{Críticas e limites}

O texto nunca menciona a Caixa de Conversão, fluxos de ouro ou conversibilidade. Câmbio aparece só para dizer que a taxa importava mais aos bancos estrangeiros (p. 149 e nota 35). O subperíodo 1906-1914 não é isolado em ponto algum: a série é anual, as regressões começam em 1907 e vão até 1930, e os anos da Caixa entram diluídos num período dominado pela guerra, pelo redesconto de 1921 e pela contração de 1923 a 1926. O modelo assume causalidade da economia para os bancos, que a própria autora admite bidirecional (p. 144). Faltam séries de juros e de câmbio, e ela lista entre as variáveis ausentes justamente as exigências do serviço da dívida externa (p. 144).

\begin{dialogo}
\item Procurar nos quatro jornais a queixa de escassez de meio circulante entre 1906 e 1914 e verificar se é atribuída à Caixa, ao Tesouro ou aos bancos. Vocabulário: falta de numerário, crise de crédito, dinheiro caro, carteira de descontos. Triner sustenta que o sistema não criava recursos financeiros (p. 138), o que deslocaria a causa da queixa para a estrutura bancária.

\item Rastrear como os jornais descrevem o Banco do Brasil no período, para testar se já lhe atribuíam papel de autoridade monetária antes do redesconto de 1921, como afirma Triner (p. 137). Vocabulário: banco de emissão, carteira cambial, banqueiro do Tesouro. Achado positivo reforça o conflito de competência entre Caixa, Tesouro e Banco do Brasil como objeto do debate.

\item Confrontar os boletins diários de movimento da Caixa com os balanços mensais dos bancos, que circulavam na mesma imprensa e são a fonte de Triner (p. 136, nota 5). Procurar articulista que ligue entrada de ouro na Caixa a expansão de depósitos ou de descontos, conexão que o artigo não examina.

\item Procurar o argumento de que a Caixa imobilizava ouro improdutivo ou drenava numerário, sobretudo em torno da custódia do ouro em 1911, e registrar que Triner não o testa. Serve para delimitar o que a literatura bancária pode arbitrar.
\end{dialogo}

\usonocapitulo{Entra na parte II como apoio de contexto, não como peça de argumento. Sustenta uma afirmação: durante a vigência da Caixa de Conversão o sistema bancário brasileiro era pequeno em nível, crescia depressa em taxa, tinha crédito pouco elástico e girava em torno de um Banco do Brasil que acumulava autoridade monetária e maior banco comercial (p. 137, p. 138, p. 149). Isso qualifica o que a Caixa podia fazer e ajuda a entender por que a disputa sobre sua competência era disputa sobre quem controlava a emissão. Serve ainda como aviso metodológico à leitura dos jornais: com depósitos subindo de 25\% para 62\% de M2 (p. 139), discutir o meio circulante apenas pelo papel-moeda emitido era discutir fração decrescente da moeda. O texto não trata da Caixa, não menciona fluxos de ouro e não separa 1906-1914, e não deve ser citado para nenhuma dessas três coisas, nem para afirmação sobre credibilidade, paridade ou regra.}
